\chapter{Evolution}
\label{chap:evolution}

This chapter introduces an evolutionary approach for improving the quality of packing solutions. 
Instead of designing a single fixed packing strategy, the algorithm evolves a population of candidate solutions and gradually improves them over time.

The main idea of evolutionary algorithms is to work with a set of candidate solutions, called \emph{individuals}, which are iteratively modified using operators inspired by natural evolution. 
In each generation, the algorithm evaluates the quality of all individuals, selects the more promising ones, and generates new candidate solutions through various evolutionary operators.
Through repeated iterations, the population is gradually refined and the algorithm searches for individuals that produce better packing solutions \cite{eiben2016}.

\section{Representation of an Individual}

An individual in the evolutionary algorithm represents the packing rules in their vector form, uniquely determining a complete packing solution.

\section{Fitness Evaluation}

The fitness of an individual corresponds directly to the fitness of the packing rules it represents. 
Individuals representing better packing rules therefore obtain higher fitness values and are more likely to persist in the population in subsequent generations.

\section{Elitist Genetic Algorithm}

The evolutionary process used in this work is based on the elitist genetic algorithm introduced by Gonçalves and Resende \cite{goncalves2013}. 
The algorithm parameters are summarized in Table~\ref{tab:ega-hyperparameters}, and the full pseudocode is given in Figure~\ref{fig:ega-pseudocode}. 
In each generation, the population is divided into two groups: \emph{elite individuals} and \emph{non-elite individuals}. 
The elite group consists of the individuals with the highest fitness values, and its size is determined by a predefined parameter.

\subsection{Selection of Parents}

After dividing the population into elite and non-elite individuals, the same number of individuals is randomly selected from both groups. 
These selected individuals form the parent population used for generating new candidate solutions.

\subsection{Uniform Crossover}

New individuals are generated by applying a \emph{uniform crossover} to pairs of elite and non-elite parents. 
This type of crossover is called \emph{uniform} because each gene of the offspring is independently chosen from one of the two parents with a given probability, rather than swapping large contiguous segments of the genome. 
The probability that a gene is inherited from the elite parent determines how strongly the elite individuals influence the new generation.

\subsection{New Population Construction}

The new population in each generation consists of three parts:

\begin{itemize}
    \item the elite individuals, which are copied directly to the next generation,
    \item the offspring generated by crossover,
    \item a set of newly generated random individuals.
\end{itemize}

This strategy combines exploitation by preserving the best individuals and exploration by introducing new random individuals.
The relative sizes of these groups are determined by the parameters of the algorithm.

\subsection{Mutation}

In addition to the original algorithm, mutation is applied to the offspring generated by crossover. 
In this approach, each gene in the packing rules vector of an offspring individual is, with a certain probability $p_M$, replaced by a new random value drawn from its valid range, an interval $[0,1]$. 

The purpose of mutation is to test whether small random modifications can help the algorithm explore new regions of the solution space. 
In combination with newly generated random individuals, mutation may help discover better packing solutions that would not be reached by crossover alone.

\subsection{Memetic}

Another extension of the original algorithm is the memetic approach.
Here, a small number of \emph{hill-climbing} iterations is applied to the elite individuals before the crossover step.

Hill climbing is a local search that applies small, mutation-like changes with a given probability.
For each iteration, a gene is randomly modified and the change is kept only if it improves fitness. 
This process is repeated for a fixed number of iterations $I_{HC}$.

We will test whether refining promising individuals this way can improve convergence and solution quality, even without adding new individuals from the rest of the population.

\begin{table}[htbp]
\centering
\caption{Elitist Genetic Algorithm Hyperparameters}
\label{tab:ega-hyperparameters}
\begin{tabular}{|l|l|}
\hline
$N$ & population size \\ \hline
$G$ & maximum number of generations \\ \hline
$p_E$ & percentage of elite individuals in the population \\ \hline
$p_C$ & probability that a gene is inherited from the elite parent during uniform crossover \\ \hline
$p_M$ & probability that a gene in a packing rules vector is mutated in offspring \\ \hline
$p_R$ & percentage of newly generated random individuals each generation \\ \hline
$I_{HC}$ & number of hill-climbing iterations applied to each elite individual \\ \hline
$p_{HC}$ & probability that a gene is mutated during each hill-climbing step \\ \hline
\end{tabular}
\end{table}

\begin{figure}[htbp]
\centering
\begin{minipage}{0.95\linewidth}
\begin{algorithmic}[1]
\State \textbf{Input:} Parameters $N$, $G$, $p_E$, $p_R$, $p_C$, $p_M$, $I_{HC}$, $p_{HC}$
\State \textbf{Output:} Final population $\mathcal{P}$ 

\State $\mathcal{P} \gets \textbf{RandomPopulation}(N)$ 
\State Evaluate($\mathcal{P}$)

\For{$g = 1$ to $G$} 
    \State $N_E = \lceil p_E \cdot N \rceil$ \Comment{Determine the number of elite individuals}
    \State $N_R = \lceil p_R \cdot N \rceil$ \Comment{Determine the number of random individuals}
    \State $N_S = N - N_E - N_R$ \Comment{Determine the number of selected parents}

    \State $\mathcal{E} \gets \text{ExtractElite}(\mathcal{P}, N_E)$ \Comment{Select the elite individuals}
    \State $\mathcal{N} = \mathcal{P} \setminus \mathcal{E}$ \Comment{Obtain the non-elite individuals}

    \State $\mathcal{E} \gets \textbf{Memetic}(\mathcal{E}, I_{HC}, p_{HC})$    \Comment{Improve elite individuals by local search}

    \State $S_E \gets \text{SelectRandom}(\mathcal{E}, N_S)$                    \Comment{Select elite parents for crossover}
    \State $S_{NE} \gets \text{SelectRandom}(\mathcal{N}, N_S)$                 \Comment{Select non-elite parents for crossover}

    \State $\mathcal{C} \gets \textbf{Crossover}(S_E, S_{NE}, p_C)$             \Comment{Generate offspring by crossover}
    \State $\mathcal{M} \gets \textbf{Mutation}(\mathcal{C}, p_M)$              \Comment{Mutate the offspring}
    \State $\mathcal{R} \gets \textbf{RandomPopulation}(N_R)$                   \Comment{Generate new random individuals}

    \State $\mathcal{P}' \gets \mathcal{E} \cup \mathcal{M} \cup \mathcal{R}$

    \State Evaluate($\mathcal{P}'$)
    \State $\mathcal{P} \gets \mathcal{P}'$
\EndFor

\State \Return $\mathcal{P}$ 
\end{algorithmic}
\end{minipage}
\caption{Elitist Genetic Algorithm Pseudocode}
\label{fig:ega-pseudocode}
\end{figure}